\chapter{Programs — Program Code}
%%% generated by the blueprint pipeline from TCSlib.Programs.Light.Compiler.ProgramCode
\section{Overview}

Declarations of the Lean module \texttt{TCSlib.Programs.Light.Compiler.ProgramCode}.

\section{Declarations}

\begin{definition}[Frame size of a program]
\lean{Light.Compiler.frameSize}
\label{Light.Compiler.frameSize}
\leanok
\uses{Light.Program, Light.Stmt.width, Light.maxOf}
For a program $P$, define its frame size to be the maximum of $3$ and the maximum width
among its statements.
\end{definition}

\begin{theorem}[Lower bound on frame size]
\lean{Light.Compiler.three_le_frameSize}
\label{Light.Compiler.three_le_frameSize}
\leanok
\uses{Light.Compiler.frameSize, Light.Program}
\difficulty{1}
For every program $P$, its frame size is at least $3$.
\end{theorem}

\begin{theorem}[Statement width bounded by frame size]
\lean{Light.Compiler.width_le_frameSize}
\label{Light.Compiler.width_le_frameSize}
\leanok
\uses{Light.Compiler.frameSize, Light.Program, Light.Stmt, Light.Stmt.width, Light.maxOf_le_iff}
\difficulty{3}
For every program $P$ and every statement $b$ in $P$, the width of $b$ is at most the
frame size of $P$.
\end{theorem}

\begin{definition}[Main call statement]
\lean{Light.mainStmt}
\label{Light.mainStmt}
\leanok
\uses{Light.Stmt}
For a procedure number $p_0$, the main statement calls procedure $p_0$ with the values
of variables $1$ and $2$, and assigns the returned value to variable $0$.
\end{definition}

\begin{definition}[Compiler tail code]
\lean{Light.Compiler.tailCode}
\label{Light.Compiler.tailCode}
\leanok
\uses{Light.Compiler.cD, Light.Compiler.cRESULT, Light.Compiler.cStack, Light.Compiler.cZERO, Light.Compiler.mainFrame, WordRAM.Instr}
Given a Boolean flag $d$ and a position $p$, the tail code first adds the value in the
main-frame stack location to the zero register and stores the result in the designated
result register. If $d$ is true, it then subtracts that stack value from zero, branches
to position $p+4$ when the difference is negative, and otherwise rejects; the final
instruction accepts. If $d$ is false, the code accepts immediately after the initial
addition.
\end{definition}

\begin{definition}[Main code position]
\lean{Light.Compiler.mainPos}
\label{Light.Compiler.mainPos}
\leanok
The main code position is the natural number $4$.
\end{definition}

\begin{definition}[Tail code position]
\lean{Light.Compiler.tailPos}
\label{Light.Compiler.tailPos}
\leanok
\uses{Light.Compiler.mainPos, Light.Stmt.size, Light.mainStmt}
For a natural number $F$, the tail code position is the main code position plus the
size, measured using parameter $F$, of the main call that calls procedure $0$ with
variables $1$ and $2$ as arguments and assigns its result to variable $0$.
\end{definition}

\begin{definition}[First body position]
\lean{Light.Compiler.firstBody}
\label{Light.Compiler.firstBody}
\leanok
\uses{Light.Compiler.tailCode, Light.Compiler.tailPos}
For a natural number $F$ and a Boolean flag $d$, the first body position is the tail
code position for $F$ plus the length of the tail code for flag $d$ at position $0$.
\end{definition}

\begin{definition}[Statement entry positions]
\lean{Light.Compiler.entries}
\label{Light.Compiler.entries}
\leanok
\uses{Light.Compiler.epilogueSize, Light.Stmt, Light.Stmt.size}
Given a natural number $F$, a list of statements, and a starting position $p$, form a
list of entry positions by recording the current position for each statement in turn.
After a statement $b$, advance the position by its size measured using $F$, plus $5$.
\end{definition}

\begin{definition}[End position]
\lean{Light.Compiler.endPos}
\label{Light.Compiler.endPos}
\leanok
\uses{Light.Compiler.epilogueSize, Light.Stmt, Light.Stmt.size}
Given a natural number $F$, a list of statements, and a starting position $p$, the end
position is obtained by advancing past each statement by its size measured using $F$,
plus $5$; for an empty list, the end position is $p$.
\end{definition}

\begin{definition}[Dispatcher position]
\lean{Light.Compiler.dispPos}
\label{Light.Compiler.dispPos}
\leanok
\uses{Light.Compiler.endPos, Light.Compiler.firstBody, Light.Compiler.frameSize, Light.Program}
For a program $P$ and a Boolean flag $d$, let $F$ be the frame size of $P$. The
dispatcher position is the end position obtained by processing the statements of $P$
starting at the first body position for $F$ and $d$.
\end{definition}

\begin{definition}[Compiler code layout]
\lean{Light.Compiler.codeLayout}
\label{Light.Compiler.codeLayout}
\leanok
\uses{Light.Compiler.CodeLayout, Light.Compiler.dispPos, Light.Compiler.entries, Light.Compiler.firstBody, Light.Compiler.frameSize, Light.Program}
For a program $P$ and a Boolean choice $d$, the code layout has frame size given by the
frame-size computation for $P$, procedure-entry positions computed from that size and
the first body selected using $d$, and a dispatcher position computed from $P$ and $d$.
\end{definition}

\begin{definition}[Compiled procedure bodies]
\lean{Light.Compiler.bodiesCode}
\label{Light.Compiler.bodiesCode}
\leanok
\uses{Light.Compiler.CodeLayout, Light.Compiler.compileStmt, Light.Compiler.epilogue, Light.Compiler.epilogueSize, Light.Stmt, Light.Stmt.size, WordRAM.Instr}
Given a code layout $G$, a list of statement bodies is compiled in order starting at
position $p$. Each body is followed by its epilogue, and the next body starts after the
current body's code and epilogue, with the epilogue length fixed at five instructions.
\end{definition}

\begin{definition}[Straight-line program header]
\lean{Light.Compiler.headStraight}
\label{Light.Compiler.headStraight}
\leanok
\uses{Light.Compiler.cNEG, Light.Compiler.cONE, Light.Compiler.cZERO, WordRAM.Instr}
The fixed header consists of three instructions: set the register for $c_{\rm ONE}$ to
$1$, set the register for $c_{\rm ZERO}$ to $0$, and set the register for $c_{\rm NEG}$ to
$-1$.
\end{definition}

\begin{definition}[Constant pool initialization]
\lean{Light.Compiler.poolCode}
\label{Light.Compiler.poolCode}
\leanok
\uses{Light.Compiler.cONE, Light.Compiler.cPool, WordRAM.Instr}
For a natural number $N$, the instruction sequence sets the pool register at index $0$
to $0$, then, for each $n\in\{0,1,\ldots,N-1\}$, sets the pool register at index $n+1$ to
the value in the pool register at index $n$ plus $1$.
\end{definition}

\begin{definition}[Zero initialization code]
\lean{Light.Compiler.zeroCode}
\label{Light.Compiler.zeroCode}
\leanok
\uses{Light.Compiler.cONE, Light.Compiler.cStack, Light.Compiler.mainFrame, WordRAM.Instr}
For a natural number $F$, the instruction sequence sets each of the $F$ stack registers
indexed by the main-frame base plus $x$, for $x\in\{0,1,\ldots,F-1\}$, to zero.
\end{definition}

\begin{definition}[Main frame initialization]
\lean{Light.Compiler.frameCode}
\label{Light.Compiler.frameCode}
\leanok
\uses{Light.Compiler.cARG1, Light.Compiler.cARG2, Light.Compiler.cFP, Light.Compiler.cPool, Light.Compiler.cStack, Light.Compiler.cZERO, Light.Compiler.mainFrame, Light.Compiler.zeroCode, WordRAM.Instr}
For a natural number $F$, the instruction sequence sets the frame pointer to the
indexed by the main-frame base plus $x$, for $x\in\{0,1,\ldots,F-1\}$, to zero.
main-frame stack registers to zero, and copies the first two argument registers into the
stack registers at offsets one and two from the main-frame base.
\end{definition}

\begin{definition}[Straight-line program startup]
\lean{Light.Compiler.startStraight}
\label{Light.Compiler.startStraight}
\leanok
\uses{Light.Compiler.frameCode, Light.Compiler.poolCode, WordRAM.Instr}
For natural numbers $F$ and $N$, the startup instruction sequence is the concatenation
of the constant-pool initialization for $N$ and the main-frame initialization for $F$.
\end{definition}

\begin{definition}[Program header code]
\lean{Light.Compiler.headCode}
\label{Light.Compiler.headCode}
\leanok
\uses{Light.Compiler.cNEG, Light.Compiler.dispPos, Light.Compiler.headStraight, Light.Program, WordRAM.Instr}
For a program $P$ and a Boolean choice $d$, the header code consists of the fixed
three-instruction setup followed by a conditional branch that tests the register holding
$-1$ and jumps to position three times the dispatcher position computed from $P$ and
$d$.
\end{definition}

\begin{definition}[Main code for a program]
\lean{Light.Compiler.mainCode}
\label{Light.Compiler.mainCode}
\leanok
\uses{Light.Compiler.codeLayout, Light.Compiler.compileStmt, Light.Compiler.frameSize, Light.Compiler.mainPos, Light.Compiler.tailCode, Light.Compiler.tailPos, Light.Program, Light.mainStmt, WordRAM.Instr}
Given a light program $P$, a procedure index $p_0$, and a Boolean flag $d$, the main
code is the concatenation of the compiled main statement with the tail code, using the
layout and instruction positions determined by $P$ and $d$.
\end{definition}

\begin{definition}[Start code for a program]
\lean{Light.Compiler.startCode}
\label{Light.Compiler.startCode}
\leanok
\uses{Light.Compiler.cNEG, Light.Compiler.dispPos, Light.Compiler.frameSize, Light.Compiler.mainPos, Light.Compiler.startStraight, Light.Program, WordRAM.Instr}
Given a light program $P$ and a Boolean flag $d$, the start code consists of the
straight-line initialization code, determined by the frame size and dispatcher position,
followed by a conditional jump to the main-code position when the designated register is
negative.
\end{definition}

\begin{definition}[Compiled program]
\lean{Light.compileProgram}
\label{Light.compileProgram}
\leanok
\uses{Light.Compiler.bodiesCode, Light.Compiler.codeLayout, Light.Compiler.dispPos, Light.Compiler.dispatcher, Light.Compiler.firstBody, Light.Compiler.frameSize, Light.Compiler.headCode, Light.Compiler.mainCode, Light.Compiler.startCode, Light.Program, WordRAM.Instr}
Given a light program $P$, a procedure index $p_0$, and a Boolean flag $d$, the compiled
program is the concatenation of the header code, main code, procedure-body code,
dispatcher, and start code, with each component determined by $P$, $p_0$, and $d$ as
appropriate.
\end{definition}

\begin{theorem}[End position is nondecreasing]
\lean{Light.Compiler.le_endPos}
\label{Light.Compiler.le_endPos}
\leanok
\uses{Light.Compiler.endPos, Light.Compiler.epilogueSize, Light.Stmt, Light.Stmt.size}
\difficulty{4}
Let $F\in\mathbb{N}$, let $bs$ be a finite list of statements, and let
$pos\in\mathbb{N}$. The end position obtained by advancing from $pos$ through the list,
adding each statement's size at frame width $F$ and an epilogue of size $5$, is at least
$pos$.
\end{theorem}

\begin{theorem}[Size of the main statement]
\lean{Light.Compiler.size_mainStmt}
\label{Light.Compiler.size_mainStmt}
\leanok
\uses{Light.Compiler.argsSize, Light.Compiler.callSize, Light.Expr.size, Light.Stmt.size, Light.mainStmt}
\difficulty{1}
For all natural numbers $p_0$ and $F$, the size of the main statement with procedure
index $p_0$, measured using frame parameter $F$, is $2F+11$.
\end{theorem}

\begin{theorem}[Width of the main statement]
\lean{Light.Compiler.width_mainStmt}
\label{Light.Compiler.width_mainStmt}
\leanok
\uses{Light.Stmt.width, Light.mainStmt}
\difficulty{1}
For every natural number $p_0$, the width of the main statement calling procedure $p_0$
with the two arguments given by variables $1$ and $2$, and storing the result in
variable $0$, is exactly $3$.
\end{theorem}

\begin{theorem}[Tail code length]
\lean{Light.Compiler.length_tailCode}
\label{Light.Compiler.length_tailCode}
\leanok
\uses{Light.Compiler.tailCode}
\difficulty{4}
For every Boolean flag $d$ and natural number $p$, the tail code has length $5$ if $d$
is true and length $2$ if $d$ is false.
\end{theorem}

\begin{theorem}[Lower bound on dispatcher position]
\lean{Light.Compiler.two_mul_frameSize_add_le_dispPos}
\label{Light.Compiler.two_mul_frameSize_add_le_dispPos}
\leanok
\uses{Light.Compiler.dispPos, Light.Compiler.firstBody, Light.Compiler.frameSize, Light.Compiler.le_endPos, Light.Compiler.length_tailCode, Light.Compiler.mainPos, Light.Compiler.size_mainStmt, Light.Compiler.tailPos, Light.Program}
\difficulty{4}
For every program $P$ and every Boolean choice of whether to use the decrementing tail
code, the dispatcher position is at least $2$ times the frame size of $P$ plus $17$.
\end{theorem}

\begin{theorem}[Length of Compiled Procedure Bodies]
\lean{Light.Compiler.length_bodiesCode}
\label{Light.Compiler.length_bodiesCode}
\leanok
\uses{Light.Compiler.CodeLayout, Light.Compiler.bodiesCode, Light.Compiler.endPos, Light.Compiler.epilogueSize, Light.Compiler.length_compileStmt, Light.Compiler.length_epilogue, Light.Stmt, Light.Stmt.size}
\difficulty{4}
Let $G$ be a code layout, let $bs$ be a finite list of statements, and let $p\in\mathbb
N$ be a starting instruction position. For each statement, compile it at its current
position and append the five-instruction epilogue; after processing a statement $b$,
advance the position by its size at frame width $G.F$ plus five. The final position
equals $p$ plus the length of the concatenated compiled code.
\end{theorem}

\begin{theorem}[Code length through main code]
\lean{Light.Compiler.length_through_mainCode}
\label{Light.Compiler.length_through_mainCode}
\leanok
\uses{Light.Compiler.codeLayout, Light.Compiler.firstBody, Light.Compiler.frameSize, Light.Compiler.headCode, Light.Compiler.headStraight, Light.Compiler.length_compileStmt, Light.Compiler.length_tailCode, Light.Compiler.mainCode, Light.Compiler.mainPos, Light.Compiler.size_mainStmt, Light.Compiler.tailPos}
\difficulty{3}
For every program $P$, main procedure index $p_0$, and Boolean flag $d$, the
concatenation of the head code and main code has length equal to the first
procedure-body position determined by the frame size of $P$ and $d$.
\end{theorem}

\begin{theorem}[Code length through procedure bodies]
\lean{Light.Compiler.length_through_bodiesCode}
\label{Light.Compiler.length_through_bodiesCode}
\leanok
\uses{Light.Compiler.bodiesCode, Light.Compiler.codeLayout, Light.Compiler.dispPos, Light.Compiler.firstBody, Light.Compiler.frameSize, Light.Compiler.headCode, Light.Compiler.length_bodiesCode, Light.Compiler.length_through_mainCode, Light.Compiler.mainCode}
\difficulty{3}
For every program $P$, main procedure index $p_0$, and Boolean flag $d$, concatenate the
head code, main code, and code for all procedure bodies in the layout determined by $P$
and $d$. The resulting list has length equal to the dispatcher position determined by
$P$ and $d$.
\end{theorem}

\begin{theorem}[Code length through the dispatcher]
\lean{Light.Compiler.length_through_dispatcher}
\label{Light.Compiler.length_through_dispatcher}
\leanok
\uses{Light.Compiler.bodiesCode, Light.Compiler.codeLayout, Light.Compiler.dispPos, Light.Compiler.dispatcher, Light.Compiler.firstBody, Light.Compiler.frameSize, Light.Compiler.headCode, Light.Compiler.length_dispatcher, Light.Compiler.length_through_bodiesCode, Light.Compiler.mainCode}
\difficulty{3}
For every program $P$, main procedure index $p_0$, and Boolean flag $d$, concatenate the
head code, main code, all procedure body code, and dispatcher code. The resulting list
has length three times the dispatcher position determined by $P$ and $d$.
\end{theorem}

\begin{theorem}[Header Location in the Compiled Program]
\lean{Light.Compiler.codeAt_headCode}
\label{Light.Compiler.codeAt_headCode}
\leanok
\uses{Light.Compiler.headCode, Light.compileProgram, WordRAM.CodeAt, WordRAM.CodeAt.left, WordRAM.codeAt_self}
\difficulty{1}
For every finite statement list $P$, natural number $p_0$, and Boolean flag $dec$, the
header code determined by $P$ and $dec$ is a prefix of the compiled program beginning at
position $0$.
\end{theorem}

\begin{theorem}[Location of the main code]
\lean{Light.Compiler.codeAt_mainCode}
\label{Light.Compiler.codeAt_mainCode}
\leanok
\uses{Light.Compiler.mainCode, Light.Compiler.mainPos, Light.compileProgram, WordRAM.CodeAt, WordRAM.CodeAt.left, WordRAM.CodeAt.right, WordRAM.codeAt_self}
\difficulty{1}
For every light program $P$, natural number $p_0$, and Boolean flag $dec$, the compiled
main code is a prefix of the compiled program beginning at position $4$.
\end{theorem}

\begin{theorem}[Location of the main statement]
\lean{Light.Compiler.codeAt_mainStmt}
\label{Light.Compiler.codeAt_mainStmt}
\leanok
\uses{Light.Compiler.codeAt_mainCode, Light.Compiler.codeLayout, Light.Compiler.compileStmt, Light.Compiler.mainPos, Light.compileProgram, Light.mainStmt, WordRAM.CodeAt, WordRAM.CodeAt.left}
\difficulty{1}
For every light program $P$, natural number $p_0$, and Boolean flag $dec$, the compiled
main statement is a prefix of the compiled program beginning at position $4$.
\end{theorem}

\begin{theorem}[Location of the terminal code]
\lean{Light.Compiler.codeAt_tailCode}
\label{Light.Compiler.codeAt_tailCode}
\leanok
\uses{Light.Compiler.codeAt_mainCode, Light.Compiler.codeLayout, Light.Compiler.frameSize, Light.Compiler.length_compileStmt, Light.Compiler.size_mainStmt, Light.Compiler.tailCode, Light.Compiler.tailPos, Light.compileProgram, WordRAM.CodeAt, WordRAM.CodeAt.cast_pos, WordRAM.CodeAt.right}
\difficulty{2}
For every light program $P$, natural number $p_0$, and Boolean flag $dec$, the terminal
code is a prefix of the compiled program beginning at the terminal-code position for the
frame size of $P$.
\end{theorem}

\begin{theorem}[Compiled Bodies Occupy Their Code Region]
\lean{Light.Compiler.codeAt_bodiesCode}
\label{Light.Compiler.codeAt_bodiesCode}
\leanok
\uses{Light.Compiler.bodiesCode, Light.Compiler.codeLayout, Light.Compiler.firstBody, Light.Compiler.frameSize, Light.Compiler.length_through_mainCode, Light.compileProgram, WordRAM.CodeAt, WordRAM.CodeAt.cast_pos, WordRAM.CodeAt.left, WordRAM.CodeAt.right, WordRAM.codeAt_self}
\difficulty{2}
For every finite statement list $P$, natural number $p_0$, and Boolean flag $dec$, the
body code for $P$, laid out from its first-body position in the layout determined by $P$
and $dec$, is a prefix of the compiled program beginning at that first-body position.
\end{theorem}

\begin{theorem}[Dispatcher Location in the Compiled Program]
\lean{Light.Compiler.codeAt_dispatcher_dispPos}
\label{Light.Compiler.codeAt_dispatcher_dispPos}
\leanok
\uses{Light.Compiler.dispPos, Light.Compiler.dispatcher, Light.Compiler.length_through_bodiesCode, Light.compileProgram, WordRAM.CodeAt, WordRAM.CodeAt.cast_pos, WordRAM.CodeAt.left, WordRAM.CodeAt.right, WordRAM.codeAt_self}
\difficulty{2}
For every finite statement list $P$, natural number $p_0$, and Boolean flag $dec$, the
dispatcher for the number of statements in $P$ is a prefix of the compiled program
beginning at the dispatcher position determined by $P$ and $dec$.
\end{theorem}

\begin{theorem}[Location of the start code]
\lean{Light.Compiler.codeAt_startCode}
\label{Light.Compiler.codeAt_startCode}
\leanok
\uses{Light.Compiler.dispPos, Light.Compiler.length_through_dispatcher, Light.Compiler.startCode, Light.compileProgram, WordRAM.CodeAt, WordRAM.CodeAt.cast_pos, WordRAM.CodeAt.right, WordRAM.codeAt_self}
\difficulty{2}
For every light program $P$, natural number $p_0$, and Boolean flag $dec$, the start
code is a prefix of the compiled program beginning at position three times the
dispatcher position.
\end{theorem}

\begin{theorem}[Code Locations for Statement Bodies]
\lean{Light.Compiler.codeAt_bodies}
\label{Light.Compiler.codeAt_bodies}
\leanok
\uses{Light.Compiler.CodeLayout, Light.Compiler.bodiesCode, Light.Compiler.compileStmt, Light.Compiler.endPos, Light.Compiler.entries, Light.Compiler.epilogue, Light.Compiler.epilogueSize, Light.Compiler.le_endPos, Light.Compiler.length_compileStmt, Light.Compiler.length_epilogue, Light.Stmt, Light.Stmt.size, WordRAM.CodeAt, WordRAM.CodeAt.cast_pos, WordRAM.CodeAt.left, WordRAM.CodeAt.right, WordRAM.Instr}
\difficulty{4}
Let $F$ be a natural number of local variables per frame, let $bs$ be a finite list of
statements, and let $pos$ be a natural-number position. Suppose the concatenation,
starting at $pos$, of the compiled code for each statement in $bs$ followed by its
epilogue is a prefix of the instruction sequence beginning at $pos$. For every index $p$
and statement $b$ occurring at index $p$ in $bs$, there is an entry position $e$ at
index $p$ in the sequence of body entry positions, such that the compiled code for $b$
followed by its epilogue is a prefix of the instruction sequence beginning at $e$, and
$e$ plus the size of $b$ (for frame size $F$) plus $5$ is at most the end position of
the bodies.
\end{theorem}
